%%%%%% -*- Coding: utf-8-unix; Mode: latex

\documentclass[polish]{projectreport}

\usepackage[utf8]{inputenc}
\usepackage{url}
\usepackage{subfigure}
\usepackage{tabularx}
\usepackage{ragged2e}
\usepackage{booktabs}
\usepackage{multirow}
\usepackage{grffile}

\usepackage[bookmarks=false]{hyperref}
\hypersetup{colorlinks,
  linkcolor=blue,
  citecolor=blue,
  urlcolor=blue}

\usepackage{indentfirst}
\usepackage{caption}
\usepackage{listings}
\usepackage[ruled,linesnumbered,lined]{algorithm2e}
\usepackage[svgnames]{xcolor}
\usepackage{inconsolata}
\usepackage{fancyhdr}

\usepackage{csquotes}
\DeclareQuoteStyle[quotes]{polish}
  {\quotedblbase}
  {\textquotedblright}
  [0.05em]
  {\quotesinglbase}
  {\fixligatures\textquoteright}
\DeclareQuoteAlias[quotes]{polish}{polish}

\usepackage[
style=numeric,
sorting=nyt,
isbn=false,
doi=true,
url=true,
backref=false,
backrefstyle=none,
maxnames=10,
giveninits=true,
abbreviate=true,
defernumbers=false,
backend=biber]{biblatex}
\addbibresource{bibliography.bib}

\lstset{
    basicstyle=\ttfamily\footnotesize,
    backgroundcolor=\color{gray!5},
    commentstyle=\it\color{Green},
    keywordstyle=\color{Red},
    stringstyle=\color{Blue},
    numberstyle=\tiny\color{Black},
    escapeinside=`',
    frame=single,
    tabsize=2,
    rulecolor=\color{black!30},
    title=\lstname,
    breaklines=true,
    breakatwhitespace=true,
    framextopmargin=2pt,
    framexbottommargin=2pt,
    extendedchars=false,
    captionpos=b,
    abovecaptionskip=5pt,
    keepspaces=true,
    numbers=left,
    numbersep=5pt,
    showspaces=false,
    showstringspaces=false,
    showtabs=false,
    tabsize=2
}

\SetAlgorithmName{\LangAlgorithm}{\LangAlgorithmRef}{\LangListOfAlgorithms}

\renewcommand{\lstlistingname}{\LangListing}
\renewcommand\lstlistlistingname{\LangListOfListings}

\newcolumntype{Y}{>{\small\centering\arraybackslash}X}

\captionsetup[figure]{skip=5pt,position=bottom}
\captionsetup[table]{skip=5pt,position=top}


%%%%%%%%%%%%%%%%%%%%%%%%%%%%%%%%%%%%%%%%%%%%%%%%%%%%%%%%%%%%%%%%%%%%%%%%%%%%%%%
\author{Dawid Czesak \and Dominik Matras}

\title{Symulacja prostego systemu spo{\l}eczno-ekonomicznego}

\date{
    {\small \texttt{dczesak@student.wszib.edu.pl}\\[1ex]
    \texttt{dmatras@student.wszib.edu.pl}\\[2ex]
    \the\year}
}

\pagestyle{fancy}
\fancyhf{}
\fancyhead[LE,RO]{\thepage}

%%%%%%%%%%%%%%%%%%%%%%%%%%%%%%%%%%%%%%%%%%%%%%%%%%%%%%%%%%%%%%%%%%%%%%%%%%%%%%%
\begin{document}

\maketitle

\thispagestyle{empty}

%%%%%%%%%%%%%%%%%%%%%%%%%%%%%%%%%%%%%%%%%%%%%%%%%%%%%%%%%%%%%%%%%%%%%%%%%%%%%%%
\section{\SectionTitleIntroduction}
\label{sec:introduction}

%%%%%%%%%%%%%%%%%%%%%%%%%%%%%%%%%%%%%%%%%%%%%%%%%%%%%%%%%%%%%%%%%%%%%%%%%%%%%%%
\subsection{Wprowadzenie}
\label{sec:intro}

Systemy spo\l{}eczno-ekonomiczne s\k{a} z\l{}o\.{z}onymi uk\l{}adami, w~kt\'orych wiele jednostek podejmuje
decyzje dotycz\k{a}ce pracy, konsumpcji, oszcz\k{e}dzania, wymiany d\'obr oraz korzystania
z~dost\k{e}pnych zasob\'ow. Zachowania pojedynczych uczestnik\'ow systemu mog\k{a} prowadzi\'{c} do
powstawania zjawisk widocznych dopiero na poziomie ca\l{}ej populacji, takich jak
nier\'owno\'sci maj\k{a}tkowe, koncentracja kapita\l{}u, zmiany poziomu dobrobytu czy r\'o\.{z}nice
w~mo\.{z}liwo\'sciach rozwoju poszczeg\'olnych grup spo\l{}ecznych.

Tradycyjne modele ekonomiczne cz\k{e}sto opisuj\k{a} gospodark\k{e} za pomoc\k{a} r\'owna\'n i~za\l{}o\.{z}e\'n
upraszczaj\k{a}cych, na przyk\l{}ad zak\l{}adaj\k{a}c racjonalno\'s\'{c} uczestnik\'ow rynku lub r\'ownowag\k{e}
systemu. W~praktyce rzeczywiste spo\l{}ecze\'nstwa i~gospodarki s\k{a} bardziej dynamiczne,
a~ich uczestnicy r\'o\.{z}ni\k{a} si\k{e} mi\k{e}dzy sob\k{a} pod wzgl\k{e}dem dochodu, maj\k{a}tku, sposobu
podejmowania decyzji, poziomu ryzyka czy dost\k{e}pu do zasob\'ow. Z~tego powodu coraz
cz\k{e}\'sciej stosuje si\k{e} symulacje komputerowe, kt\'ore pozwalaj\k{a} obserwowa\'{c}, jak proste
regu\l{}y zachowania pojedynczych jednostek wp\l{}ywaj\k{a} na dzia\l{}anie ca\l{}ego systemu.

Jednym z~popularnych podej\'s\'{c} do badania takich zjawisk jest modelowanie agentowe,
czyli \emph{Agent-Based Modeling} (ABM). W~tego typu modelach system sk\l{}ada si\k{e}
z~wielu autonomicznych agent\'ow, kt\'orzy posiadaj\k{a} okre\'slone cechy, podejmuj\k{a} decyzje
oraz wchodz\k{a} ze sob\k{a} w~interakcje. Dzi\k{e}ki temu mo\.{z}na analizowa\'{c}, jak lokalne dzia\l{}ania
pojedynczych uczestnik\'ow prowadz\k{a} do globalnych efekt\'ow, takich jak wzrost nier\'owno\'sci,
stabilizacja systemu lub jego destabilizacja.

%%%%%%%%%%%%%%%%%%%%%%%%%%%%%%%%%%%%%%%%%%%%%%%%%%%%%%%%%%%%%%%%%%%%%%%%%%%%%%%
\subsection{Modelowanie agentowe w systemach spo\l{}eczno-ekonomicznych}
\label{sec:abm}

Modelowanie agentowe jest szczeg\'olnie przydatne w~analizie system\'ow
spo\l{}eczno-ekonomicznych, poniewa\.{z} umo\.{z}liwia uwzgl\k{e}dnienie r\'o\.{z}norodno\'sci uczestnik\'ow
systemu. Ka\.{z}dy agent mo\.{z}e mie\'{c} inne zasoby finansowe, inne potrzeby, inne strategie
dzia\l{}ania oraz inn\k{a} sk\l{}onno\'s\'{c} do oszcz\k{e}dzania lub wydawania pieni\k{e}dzy. W~przeciwie\'nstwie
do modeli opartych wy\l{}\k{a}cznie na warto\'sciach \'srednich, modele agentowe pozwalaj\k{a}
obserwowa\'{c} zachowanie ca\l{}ej populacji oraz pojedynczych jednostek.

W~kontekcie prostego systemu spo\l{}eczno-ekonomicznego agentami mog\k{a} by\'{c} na przyk\l{}ad
osoby, gospodarstwa domowe, firmy lub instytucje. Agenci mog\k{a} posiada\'{c} okre\'slony
maj\k{a}tek pocz\k{a}tkowy, otrzymywa\'{c} doch\'od, ponosi\'{c} koszty \.{z}ycia, handlowa\'{c} z~innymi
agentami, p\l{}aci\'{c} podatki lub korzysta\'{c} z~pomocy spo\l{}ecznej. Zmiana jednego parametru,
na przyk\l{}ad wysoko\'sci podatku albo sposobu redystrybucji \'srodk\'ow, mo\.{z}e wp\l{}ywa\'{c} na
ko\'ncowy poziom nier\'owno\'sci w~ca\l{}ej populacji.

W~badaniach naukowych modele agentowe s\k{a} wykorzystywane mi\k{e}dzy innymi do analizy
rynk\'ow, przep\l{}ywu pieni\k{e}dzy, nier\'owno\'sci dochodowych, wp\l{}ywu polityki spo\l{}ecznej,
zachowa\'n konsument\'ow oraz proces\'ow gospodarczych. Przegl\k{a}d bada\'n z~2024~roku dotycz\k{a}cy
zastosowania modelowania agentowego w~gospodarce cyrkularnej wskazuje, \.{z}e metoda ta
pozwala analizowa\'{c} zachowania producent\'ow, konsument\'ow oraz procesy dyfuzji innowacji
w~systemach ekonomicznych~\cite{mudd2024}.

%%%%%%%%%%%%%%%%%%%%%%%%%%%%%%%%%%%%%%%%%%%%%%%%%%%%%%%%%%%%%%%%%%%%%%%%%%%%%%%
\subsection{Przegl\k{a}d aktualnego stanu bada\'n}
\label{sec:related}

W~ostatnich latach modelowanie agentowe jest coraz cz\k{e}\'sciej stosowane do badania
nier\'owno\'sci spo\l{}ecznych i~ekonomicznych. Punktem wyj\'scia dla orientacji w~tym obszarze
mo\.{z}e by\'{c} systematyczny przegl\k{a}d Mudda i~wsp\'o\l{}autor\'ow~\cite{mudd2024}, kt\'orzy skatalogowali
modele symulacyjne i~koncepcyjne dotycz\k{a}ce nier\'owno\'sci spo\l{}eczno-ekonomicznych
w~zdrowiu i~zachowaniach zdrowotnych. Przegl\k{a}d ten ujawni\l{}, \.{z}e wi\k{e}kszo\'s\'{c} istniej\k{a}cych
modeli koncentruje si\k{e} na po\'srednich determinantach zdrowia, pomijaj\k{a}c strukturalne
czynniki systemowe. Niniejsza praca przyjmuje odmienne podej\'scie -- zamiast mapowa\'{c}
istniej\k{a}ce modele, budujemy w\l{}asn\k{a} symulacj\k{e} od podstaw, co pozwala pe\l{}niej
kontrolowa\'{c} za\l{}o\.{z}enia i~parametry systemu.

Bezpo\'srednio zwi\k{a}zane z~tematyk\k{a} niniejszej pracy s\k{a} badania nad mechanizmami powstawania
nier\'owno\'sci dochodowych i~maj\k{a}tkowych. Palagi i~wsp\'o\l{}autorzy~\cite{palagi2023} zbudowali
model agentowy badaj\k{a}cy jednocze\'snie dynamik\k{e} wzrostu i~podzia\l{}u dochod\'ow, uwzgl\k{e}dniaj\k{a}c
ograniczenia kredytowe i~nier\'owny dost\k{e}p do inwestycji. Co istotne, zamiast tradycyjnego
efektu \emph{trickle-down}, zaobserwowali oni odwrotny mechanizm -- redystrybucja
w~stron\k{e} biedniejszych gospodarstw zwi\k{e}ksza\l{}a popyt i~korzystnie wp\l{}ywa\l{}a na doch\'ody
ca\l{}ej populacji. W~odniesieniu do naszego projektu, powy\.{z}sza praca potwierdza, \.{z}e ju\.{z}
stosunkowo prosty model agentowy mo\.{z}e odtworzy\'{c} nieliniowe zale\.{z}no\'sci mi\k{e}dzy
redystrybucj\k{a} a~wzrostem -- zjawisko trudne do uchwycenia w~modelach analitycznych.

Innym reprezentatywnym podej\'sciem jest simulator TaxAI~\cite{mi2024taxai}, kt\'ory stawia
sobie znacznie bardziej ambitny cel: optymalizacj\k{e} polityki podatkowej przy u\.{z}yciu
wieloagentowego uczenia ze wzmocnieniem na populacji nawet 10\,000 agent\'ow. O~ile
TaxAI dostarcza zaawansowanego narz\k{e}dzia dla decydent\'ow, jego z\l{}o\.{z}ono\'s\'{c}
utrudnia interpretacj\k{e} podstawowych mechanizm\'ow rz\k{a}dz\k{a}cych nierównowagi\.{a} w~systemie.
Nasza symulacja \'swiadomie rezygnuje z~mechanizm\'ow uczenia maszynowego na rzecz
przejrzysto\'sci -- celem jest zrozumienie, nie optymalizacja.

Najbli\.{z}szy tematycznie jest model Villafa\~{n}e i~wsp\'o\l{}autor\'ow~\cite{villafane2025}, oparty
na schemacie Yard-Sale, w~kt\'orym agenci losowo wymieniaj\k{a} mi\k{e}dzy sob\k{a} cz\k{e}\'s\'{c} maj\k{a}tku.
Badania pokaza\l{}y, \.{z}e bez mechanizm\'ow koryguj\k{a}cych nawet tak proste regu\l{}y transakcyjne
prowadz\k{a} nieuchronnie do skrajnej koncentracji zasob\'ow. Nasz model inspiruje si\k{e}
tym podej\'sciem, rozszerzaj\k{a}c je o~doch\'od, koszty \.{z}ycia i~proste formy redystrybucji,
co pozwala zbada\'{c} warunki brzegowe powstawania i~ograniczania nier\'owno\'sci.

Szerszy kontekst zastosowa\'n ABM w~instytucjach decyzyjnych zarysowuj\k{a} Borsos
i~wsp\'o\l{}autorzy~\cite{borsos2025} w~przegl\k{a}dzie obejmuj\k{a}cym dwie dekady bada\'n bank\'ow
centralnych. Zestawienie to pokazuje, \.{z}e modele agentowe przesz\l{}y d\l{}ug\k{a} drog\k{e} od narz\k{e}dzi
akademickich do realnych narz\k{e}dzi analitycznych wspieraj\k{a}cych polityk\k{e} makroostro\.{z}no\'sciow\k{a}.
Stanowi to dodatkowe uzasadnienie dla podej\'scia przyj\k{e}tego w~niniejszej pracy --
nawet uproszczony model agentowy mo\.{z}e pe\l{}ni\'{c} rol\k{e} u\.{z}ytecznego narz\k{e}dzia do analizy
mechanizm\'ow spo\l{}eczno-ekonomicznych.

%%%%%%%%%%%%%%%%%%%%%%%%%%%%%%%%%%%%%%%%%%%%%%%%%%%%%%%%%%%%%%%%%%%%%%%%%%%%%%%
\subsection{Cel pracy}
\label{sec:goal}

Celem pracy jest opracowanie prostej symulacji systemu spo\l{}eczno-ekonomicznego opartej
na modelu agentowym. W~proponowanym modelu populacja sk\l{}ada si\k{e} z~grupy agent\'ow
reprezentuj\k{a}cych uczestnik\'ow systemu ekonomicznego. Ka\.{z}dy agent posiada okre\'slony
poziom maj\k{a}tku, mo\.{z}e uzyskiwa\'{c} doch\'od, ponosi\'{c} koszty \.{z}ycia oraz wchodzi\'{c}
w~interakcje ekonomiczne z~innymi agentami.

Symulacja ma umo\.{z}liwi\'{c} obserwacj\k{e} zmian zachodz\k{a}cych w~czasie, w~szczeg\'olno\'sci:
\begin{itemize}
  \item zmian poziomu maj\k{a}tku poszczeg\'olnych agent\'ow,
  \item powstawania nier\'owno\'sci ekonomicznych,
  \item wp\l{}ywu prostych regu\l{} wymiany i~redystrybucji na struktur\k{e} maj\k{a}tkow\k{a} populacji,
  \item zale\.{z}no\'sci mi\k{e}dzy decyzjami pojedynczych agent\'ow a~zachowaniem ca\l{}ego systemu.
\end{itemize}

Praca ma charakter eksperymentalny. Jej celem nie jest dok\l{}adne odwzorowanie
rzeczywistej gospodarki, lecz pokazanie, w~jaki spos\'ob za pomoc\k{a} prostego modelu
komputerowego mo\.{z}na analizowa\'{c} podstawowe mechanizmy spo\l{}eczno-ekonomiczne.

%%%%%%%%%%%%%%%%%%%%%%%%%%%%%%%%%%%%%%%%%%%%%%%%%%%%%%%%%%%%%%%%%%%%%%%%%%%%%%%
\subsection{Zakres dalszej cz\k{e}\'sci sprawozdania}
\label{sec:scope}

W~dalszej cz\k{e}\'sci sprawozdania zostanie przedstawione sformu\l{}owanie problemu badawczego
oraz opis zaproponowanego rozwi\k{a}zania. Nast\k{e}pnie opisany zostanie model symulacyjny,
jego za\l{}o\.{z}enia, parametry oraz zastosowane algorytmy. Kolejna cz\k{e}\'s\'{c} b\k{e}dzie zawiera\l{}a
wyniki eksperyment\'ow symulacyjnych, ich analiz\k{e} oraz interpretacj\k{e}. Na ko\'ncu pracy
zostan\k{a} przedstawione wnioski oraz mo\.{z}liwe kierunki dalszego rozwoju modelu.

% %%%%%%%%%%%%%%%%%%%%%%%%%%%%%%%%%%%%%%%%%%%%%%%%%%%%%%%%%%%%%%%%%%%%%%%%%%%%%%%
% \section{\SectionTitleProblemSolution}
% \label{sec:solution}

% \emph{Sformu\l{}owanie problemu badawczego. Propozycja rozwi\k{a}zania: szczeg\'o\l{}owy opis
% algorytm\'ow/systemu/frameworka.}

% %%%%%%%%%%%%%%%%%%%%%%%%%%%%%%%%%%%%%%%%%%%%%%%%%%%%%%%%%%%%%%%%%%%%%%%%%%%%%%%
% \section{\SectionTitleExperiments}
% \label{sec:experiments}

% \emph{Wyniki bada\'n eksperymentalnych wraz ze szczeg\'o\l{}owymi analizami i~komentarzami.}

% %%%%%%%%%%%%%%%%%%%%%%%%%%%%%%%%%%%%%%%%%%%%%%%%%%%%%%%%%%%%%%%%%%%%%%%%%%%%%%%
% \section{\SectionTitleSummary}
% \label{sec:conclusions}

% \emph{Podsumowanie wynik\'ow bada\'n, wnioski oraz plany dalszych prac.}

%%%%%%%%%%%%%%%%%%%%%%%%%%%%%%%%%%%%%%%%%%%%%%%%%%%%%%%%%%%%%%%%%%%%%%%%%%%%%%%
\printbibliography

\end{document}